% Paper 2 (ReSIDS v2) -- dataset audit. Impersonal style.
% Every number below comes from a released script and its report:
%   scripts/inspect_edge_iiotset.py   -> results/edge_iiotset_inspection.txt
%   scripts/inspect_uav_nidd.py       -> results/uav_nidd_inspection.txt
%   scripts/list_uav_nidd_figshare.py -> results/uav_nidd_figshare_listing.txt
%   scripts/inspect_pdnet_ids26.py    -> results/pdnet_ids26_inspection.txt
%   scripts/inspect_uav_hassler.py    -> results/uav_hassler_inspection.txt
%   scripts/inspect_ereno.py          -> results/ereno_audit.txt (+ ereno_attribution.txt)
%   scripts/inspect_ciciot2023.py     -> results/ciciot2023_audit.txt (10 of 169 parts)
% Citation keys are in docs/refs_v2.bib (section "Dataset audit").
%   scripts/inspect_cicids2017c.py    -> results/cicids2017c_audit.txt
%   scripts/inspect_iomt2024.py       -> results/iomt2024_audit.txt (Kaggle mirror; checksums in results/iomt2024_md5sums.txt)
%   scripts/inspect_cic_iiot2025.py   -> results/cic_iiot2025_audit.txt (Kaggle mirror; checksums in results/cic_iiot2025_md5sums.txt)
%   scripts/inspect_cps_survey4.py    -> results/{morris1,hai2007,westermo,powerduck}_audit.txt
%                                        (+ results/cps_survey4_followup.txt)
%   scripts/inspect_hai2305.py        -> results/hai2305_audit.txt (SHA-256 checked against GitHub LFS)

\subsection{Dataset audit}
\label{sec:dataset-audit}

The evaluation of Section~\ref{sec:datasets} requires properties that accuracy
benchmarks seldom exercise: rows aligned with their header, labels that describe attack
behaviour rather than the capture session, a time order from which windows can be built,
source addresses for attribution, and enough benign traffic to measure precision at
operational prevalence. Ten public datasets from cyber-physical and IoT settings were
considered as additional domains: two industrial IoT, two medical, two aerial, two power
system and two industrial control system datasets. The last four were shortlisted from a
recent survey of CPS intrusion datasets~\cite{quincozes2025cpsdatasets}, whose other
smart-grid and ICS entries were excluded on their published description: EPIC has no
per-sample labels, Biswas et al.\ and ICS-ADD are unlabelled packet captures, the Sahu
dataset is not public, and ICS-Flow and Electra carry no physical data. Before use, each of
the ten was audited with the same protocol, and none could be retained. The three datasets used in the evaluation,
ERENO, CICIoT2023 and the corrected CICIDS2017, went through the same protocol. Table~\ref{tab:dataset-audit}
summarises the outcome; the scripts and their reports are released with the artifact.

\paragraph{Protocol.}
Six checks were applied to every release.
(i)~\emph{Structure}: fields per line, embedded header lines, and agreement between the
kind of value a column holds (epoch time, radio frequency, signal power, IPv4 address)
and the name in its header.
(ii)~\emph{Labels}: exact duplicates, duplicates once identifiers are removed, feature
vectors carrying more than one label, and identical files filed under different labels.
(iii)~\emph{Time}: presence and completeness of timestamps, and whether classes are
interleaved or recorded in separate sessions.
(iv)~\emph{Attribution}: whether source addresses exist and point to the emitter.
(v)~\emph{Prevalence}: class balance of the release and benign volume.
(vi)~\emph{Shortcut probe}: a gradient-boosted classifier trained on deduplicated rows
with and without the columns that locate a row in a capture (timestamps, frame and
sequence numbers, checksums, ports, addresses), on a random split and, where the release
provides one, on an independent test file. Three outcomes indicate that labels are
predictable from the capture rather than from the attack: separability carried by one
such column, a collapse once they are removed, or separability that holds on a random
split but not on independent data.

\begin{table}[t]
\centering
\caption{Outcome of the dataset audit. Checks: S structure, L labels, T time,
A attribution, P prevalence, X shortcut probe.}
\label{tab:dataset-audit}
\small
\begin{tabular}{@{}lllll@{}}
\toprule
Dataset & Domain & Failed checks & Main finding & Use \\
\midrule
ERENO~\cite{quincozes2023ereno} & IEC 61850 substation & A &
  attacks impersonate the single publisher & retained \\
CICIoT2023~\cite{neto2023ciciot2023} & IoT network & T, A, P &
  no time or address; IAT acts as a capture clock & retained \\
CICIDS2017, corrected~\cite{engelen2021troubleshooting} & enterprise network & L, X &
  scans duplicated across two labels; window-size fingerprint & retained \\
\midrule
Edge-IIoTset~\cite{ferrag2022edgeiiotset} & industrial IoT & S, L, T, A, X &
  two classes shifted by one column & rejected \\
UAV-NIDD~\cite{hadi2025uavnidd} & UAV network & S, L, T, A, P, X &
  captures reused across scenarios and labels & rejected \\
PDNet-IDS26~\cite{lui2026pdnet} & medical IoT & P, X &
  class identified by the TCP connection & rejected \\
Hassler et al.~\cite{hassler2024uav} & UAV, cyber-physical & S, T, X &
  two attacks without benign data in their schema & rejected \\
CIC-BCCC-NRC-IoMT-2024~\cite{dadkhah2024ciciomt} & medical IoT & L, T, A, P, X &
  ICMP, UDP and ARP classes hold only TCP flows & rejected \\
CIC IIoT 2025~\cite{firouzi2025datasense} & industrial IoT, sensors & L, T, A, X &
  empty windows labelled attack; sensors carry no attack signal & rejected \\
Morris-1~\cite{pan2015hybrid} & power system (PMU) & T, A, P &
  1\% random sample: no time order; 71\% attacks & rejected \\
PowerDuck~\cite{zemanek2022powerduck} & IEC 61850 substation & A, P, X &
  floods are 92.5\% of packets; no transfer across recordings & rejected \\
Westermo~\cite{westermo2023dataset} & ICS network & L &
  timing labels mark flows not involving the attacker; no process data & not used \\
HAI 23.05~\cite{shin2020hai} & ICS process (HIL) & A, X &
  sound labels; supervised detection fails on unseen attacks & not used \\
\bottomrule
\end{tabular}
\end{table}

\paragraph{Edge-IIoTset.}
The selected release of Edge-IIoTset~\cite{ferrag2022edgeiiotset} holds 2{,}219{,}201
packets and 14 attacks. Every row of two classes, DDoS-UDP (121{,}568) and MITM
(1{,}214), is shifted by one column: the timestamp field holds an IP address or a
protocol value, so the features of these classes describe other fields, and the classes
become separable by an export artefact. The reduced release distributed for machine
learning shows the same shift in its MITM rows. The timestamp also lost its month and
day when the original date was split at its comma; captures collected between November
2021 and January 2022 therefore overlap in time of day, and the file is ordered in one
block per class. Labels are assigned per capture and applied to every packet: once the
identifiers recommended by the authors are removed, 13.95\% of the rows are duplicates
and 1{,}513 feature vectors carry more than one label. The probe reaches a macro-F1 of
0.51 on the cleaned features and 0.15 once per-packet arbitrary values (sequence and
acknowledgement numbers, checksums, stream identifiers) are also removed, so
separability rests on values that identify captures rather than attacks. Source
addresses cannot support attribution: benign traffic comes from two hosts, and about
half of the attack packets carry as source the address of a host that also originates
41\% of the benign traffic.

\paragraph{UAV-NIDD.}
UAV-NIDD~\cite{hadi2025uavnidd,hadi2025uavniddData} provides labelled files for three
compromise scenarios (UAV, access point, ground control station) and the raw captures.
In the UAV scenario, 24{,}732 of 885{,}384 lines (2.8\%) are truncated to one to five
fields, and several columns hold another field's values: the relative-time column holds
epoch timestamps and the source-address column holds protocol numbers. The
access-point file names its label column ``Normal'' and shows the same kind of shift.
The timestamps place the classes in separate sessions years apart: benign traffic of the
UAV scenario dates from September 2017 and most attacks from March 2024. On the
access-point file, the probe separates all classes perfectly (macro-F1 1.00) even
without time, position and identity columns. In the ground-control-station file, the
50{,}135 reconnaissance flows are a row-by-row copy of the 50{,}135 scanning flows, and
76\% of the rows are duplicates; the file published later as an update has the same
size and the same per-class counts. The raw release explains these observations.
Checksums show that the three scenarios reuse the same captures: the brute-force and
fake-landing captures appear in all three, and one scanning capture is filed as scanning
in two scenarios and as reconnaissance in a third. Identical files are also filed under
different attacks, de-authentication and MITM, and fake landing and replay. Finally,
some denial-of-service captures are named after UDP ports 5554 and 5556, the
navigation-data and command ports of the Parrot AR.Drone protocol, a platform that is
not part of the described testbed.

\paragraph{PDNet-IDS26.}
PDNet-IDS26~\cite{lui2026pdnet} fuses MQTT traffic with Parkinson's telemonitoring
records and is the only audited release with clinical payloads. The public file is
undersampled to 25\% per class (5{,}388 rows each), so precision at operational
prevalence cannot be measured, as its authors acknowledge. Each class arrives on a
single subscriber TCP connection (two for replay). The destination port alone separates
FDI from benign traffic with 100\% accuracy, and the probe reaches a macro-F1 of 1.00 on
every class even after patient identity and sequence numbers are removed. The F1-score
of 0.77 reported for FDI by the authors indicates that the published file retains a
session identifier that their own evaluation did not use. Denial-of-service rows carry
an all-zero clinical payload, and the replay class covers 9 of the 42 patients.

\paragraph{Hassler et al.}
The cyber-physical UAV dataset of Hassler et al.~\cite{hassler2024uav,hassler2024uavData}
stacks ten blocks in one CSV, one network and one physical block per class, under five
schemas. Evil twin and FDI share a network schema that contains no benign rows, and
their physical blocks use two further schemas, so neither attack has benign data in its
own feature space. The FDI physical block has no timestamp, and the evil-twin physical
clock is offset by five hours from its network clock, consistent with a time-zone error.
Benign traffic, de-authentication DoS and replay share schemas and can be aligned in
time, with offsets of 16 to 48~s. Yet the probe on their network features falls from a
macro-F1 of 0.99 to 0.65 once timestamps and sequence numbers are removed, because the
classes were recorded in separate flight sessions.

\paragraph{CIC-BCCC-NRC-IoMT-2024.}
This release re-extracts CICIoMT2024~\cite{dadkhah2024ciciomt}, captured from 40 medical
IoT devices, into CICFlowMeter flows with addresses and timestamps, within the
TabularIoTAttacks-2024 collection~\cite{sasi2024tabular}. Its structure is sound: one
header, 85 fields per line, and well-formed addresses, ports and timestamps. Its labels
describe capture sessions, not flows. All 3{,}385{,}313 flows are TCP, including those of
the ICMP-flood, UDP-flood and ARP-spoofing classes; in those classes most flows go to
MQTT (port 1883) and HTTPS (port 443), that is, ordinary device traffic recorded while
the attack ran. Each class was recorded in its own session between 28~August and
12~September 2023, and benign traffic only on 13 and 14~September, so no benign flow
falls inside any attack window. Benign flows are 0.96\% of the release, DDoS flows
0.15\% and ARP spoofing 0.03\%. The probe falls from a macro-F1 of 0.881 on a random
split with identifiers to 0.769 on a time-ordered split without them, and DDoS and ARP
spoofing reach only 0.39 and 0.37, consistent with classes whose flows do not contain
the attack. The audit used a third-party mirror of the release (checksums released with the
scripts).

\paragraph{CIC IIoT 2025.}
CIC IIoT 2025 (DataSense)~\cite{firouzi2025datasense} is the only candidate that joins
network traffic with industrial sensor readings. Its processed release describes each of
38 devices per time window, from 1 to 10~s, with sensor-log and network aggregates. The
audit used the 1-s and 10-s windows. Their structure is sound, but three checks fail.
Labels: half of the benign windows (50.6\%) and 10.7\% of the attack windows contain
neither traffic nor sensor readings, so 9{,}748 empty windows are labelled as attacks.
Time: the benign file covers one hour on 9~September 2025, while the attacks were recorded
in separate sessions between January and June 2025, so no benign window falls inside any
attack span. The probe falls from a macro-F1 of 0.968 on a random split to 0.726 on a
time-ordered split, with brute force, DoS and DDoS between 0.34 and 0.59. The release is
also inconsistent across window sizes: the ten-second attack file has 5.5 times fewer
windows than the one-second file, where 10 would be expected, so the two resolutions do
not tile the same attack time. Above all, the sensor data carry no attack signal. Trained
on sensor features alone, the probe reaches a macro-F1 of 0.17, against 0.97 with network
features. Moreover, sensors report in 99\% of their benign windows but in about half of
their attack windows, a difference between sessions rather than an effect of the attacks.
The physical part of this dataset therefore cannot support a cyber-physical evaluation.

\paragraph{Power system and industrial control datasets.}
Morris-1~\cite{pan2015hybrid} joins synchrophasor measurements with relay, control-panel
and Snort logs for natural faults and attacks on a transmission system. Its structure is
sound: 78{,}377 rows, no conflicting labels and negligible duplication, although 10{,}906
impedance values are infinite. Each of its 15 sets, however, is a 1\% random sample of the
scenario recordings and carries no timestamp, so no time-ordered window can be built.
There is no address column; 71\% of the rows are attacks; and only the binary labelling
remains online. Trained on ten sets and tested on the other five, the probe reaches a
macro-F1 of 0.60, with physical measurements carrying almost all of the signal.

PowerDuck~\cite{zemanek2022powerduck} records GOOSE traffic in a physical substation
testbed and labels every attacker packet; its labels agree with the published attack list
for 99.97\% of packets. Two flooding recordings, however, make up 92.5\% of all packets
(5.6\% malicious without them), and 92\% of the packet vectors are duplicates. As in
ERENO, every malicious packet carries the legitimate publisher's address. Most
importantly, separability does not carry across recordings. On deduplicated data the
probe reaches a macro-F1 of 0.99 on a random split, but when a recording is held out,
while other recordings of the same attack family remain in training, the F1-score of the
malicious class falls to 0.00 for replay, 0.07 for insertion and 0.30 for suppression.

Westermo~\cite{westermo2023dataset} records 90 minutes of traffic in an industrial
network, with scheduled events interleaved in one continuous capture and source addresses
that identify the attacker. It provides two labellings. The timing labelling marks every
flow during an event as malicious, and 67\% of the flows it marks do not involve the
attacker. The attacker-involvement labelling is sound and needs no session shortcut: the
probe scores the same on random and time-ordered splits (macro-F1 0.61 and 0.60). The
flow files carry no process data, however, and two of the six events are
misconfigurations invisible in flow statistics (F1-score near zero). The dataset therefore
adds no property beyond the corrected CICIDS2017.

HAI 23.05~\cite{shin2020hai}, the latest release of the HIL-based augmented ICS dataset
(steam-turbine and pumped-storage generation emulated in the loop), is the soundest of the
candidates. It holds 86 process variables at 1~Hz with no gaps, 249 hours of normal
operation for training, and 52 attacks interleaved with normal operation (4.0\% of the
test time). Its binary labels match, segment by segment, the start time and duration of
the 52 attacks listed in the technical manual. The only defect is that one label file
truncates its timestamps to the minute, so it aligns with the data by row order only.
Joining the labels with the manual yields a per-attack label set by target control loop.
Supervised detection, however, does not transfer beyond the attacks it has seen. The
probe reaches a macro-F1 of 0.99 on a random split, but 0.53 when trained on one test file
and tested on the other, 0.27 on a time-ordered split, and 0.07 per control loop on
held-out attack instances. Each attack instance is distinct, normal operation drifts
across days, and consecutive seconds are strongly correlated, which inflates any random
split. The earlier release (HAI 20.07) behaves better (0.86 across files, 0.92
time-ordered), but its labels name only the attacked process. HAI is built for one-class
anomaly detection trained on normal operation. Using it would require one-class
specialists, a change of detector design outside the scope of this paper, and it has no
network traffic for attribution.

\paragraph{Implications.}
These findings do not make the releases useless for every purpose. They do prevent an
evaluation built on time-ordered windows, source attribution and realistic prevalence,
and results reported on these releases should be read with them in mind. Two causes
recur. The first is per-session labelling: each class is recorded in its own capture, so
time, session and sequence fields identify the class. The second is export errors:
shifted columns, truncated lines, merged schemas and duplicated files. A pre-release
check of field counts, of value kinds per column, of checksums across files, and a
probe with identifiers removed would expose all of them.

\paragraph{Retained datasets.}
ERENO~\cite{quincozes2023ereno} passes the structure, label and time checks. Every line
has 59 fields; no feature vector carries two labels; duplicates are below 1\%; no test
vector reappears in the training file; and attacks are interleaved with benign traffic
over the whole simulation, with 97 to 100\% of every attack class inside the benign time
span. The probe stays near-perfect on the authors' independent test file (macro-F1 0.994)
after absolute time and the SqNum/StNum counters are removed, so separability is carried
by message content rather than by capture position. It also reflects that the attacks
are generated by deterministic rules, a limitation of synthetic data. ERENO fails the
attribution check. All benign messages and 99.9 to 100\% of the attack messages carry
the same publisher identity (one source MAC address and one \texttt{gocbRef}), because
the attacks impersonate the legitimate publisher. Source attribution in this domain would
therefore point at the victim, and isolation is exercised only in synthetic scenarios
(Section~\ref{sec:api-events}).

CICIoT2023~\cite{neto2023ciciot2023} was audited on 10 of the 169 files of its CSV
release (2{,}366{,}956 rows). It passes the structure and label checks: one header, 47
fields per line, no missing or infinite values, no duplicate rows, no feature vector with
two labels, and rows in shuffled order. It has neither timestamps nor addresses, so the
local tick is count-based and attribution is none (Section~\ref{sec:datasets}). The
column named IAT is not an inter-arrival time. Its values form narrow bands per class,
around $8.3\times10^{7}$ for denial-of-service and Mirai traffic and $1.67\times10^{8}$
for benign, reconnaissance, spoofing and web traffic, as a capture clock would. Removing
it lowers the probe's F1-score for DoS from 1.00 to 0.73, so the DoS/DDoS distinction
partly rests on it, and IAT is therefore excluded from the features. Benign traffic is
2.4\% of the sample, and the web and brute-force categories together 0.08\%: benign
traffic must be rebalanced to a stated prevalence, and those two categories yield wide
confidence intervals.

The corrected CICIDS2017~\cite{engelen2021troubleshooting,liu2022error} (2{,}099{,}976
flows over five days) passes the structure, time, attribution and prevalence checks.
Every line has 91 fields; addresses, ports and timestamps are well formed; and attacks
run concurrently with benign traffic, which makes up 15 to 99.8\% of the flows inside
each attack's time window. Attack flows come from the external attacker's gateway
(172.16.0.1) or from the internal hosts compromised by the botnet and the infiltration,
so source attribution points at the emitter. Monday holds eight hours of benign-only
traffic (371{,}624 flows). Two findings set conditions of use. First, 126 feature
vectors carry more than one label, covering 91{,}253 rows that are almost all port scans
labelled both ``Portscan'' and ``Infiltration -- Portscan''. Scans launched from the
infiltrated host are indistinguishable from those launched from outside, so the two
labels are merged. Second, the probe falls from a macro-F1 of 0.991 on a random split to
0.932 on a time-ordered split, and to 0.908 once the initial TCP window sizes are also
removed. In that last step the port-scan F1-score drops from 0.86 to 0.62. The initial
window acts as a fingerprint of the scanning host's operating system rather than of the
attack, and it is therefore excluded from the features.
Attempted flows (11{,}979), whose attacks delivered no payload, are labelled benign; the
web and infiltration categories keep only 104 and 36 flows.
